\chapter{Présentation de l'organisme d'accueil}
\clearpage

\section{Introduction}

Ce chapitre présente l'organisme d'accueil, Hightech Payment Systems (HPS). Il introduit son positionnement, son historique, ses activités principales, son organisation interne ainsi que la direction et l'équipe au sein desquelles le stage a été réalisé. Une attention particulière est accordée à la Direction Software Factory et à l'équipe Software Security, dont les missions sont directement liées à la sécurité logicielle, au cycle de vie du développement et à la conformité PCI SSF.

\section{Présentation de HPS}

Hightech Payment Systems, généralement désignée par l'acronyme \textbf{HPS}, est une entreprise multinationale marocaine spécialisée dans les solutions de paiement électronique. Elle s'adresse principalement aux institutions financières, aux processeurs de paiement ainsi qu'aux switches nationaux et régionaux. Son activité se situe au croisement de la monétique, de l'ingénierie logicielle et de la cybersécurité applicative.

Fondée à Casablanca, HPS s'est progressivement imposée comme un acteur important du secteur des paiements numériques. Son offre repose notamment sur la suite logicielle \textbf{PowerCARD}, une plateforme destinée à couvrir plusieurs besoins liés à l'émission de cartes, à l'acquisition, au traitement des transactions, au switch, à la lutte contre la fraude et à la gestion de services de paiement via différents canaux tels que les guichets automatiques, les terminaux de paiement, le web et le mobile.

Les solutions développées par HPS sont utilisées dans des environnements où la disponibilité, la performance, la confidentialité et l'intégrité des données constituent des exigences majeures. Cette réalité explique l'importance accordée par l'entreprise aux référentiels de sécurité, à la conformité réglementaire et aux pratiques de développement sécurisé, en particulier dans le cadre du \textit{PCI Secure Software Framework} (PCI SSF).

\begin{figure}[h!]
    \centering
    % TODO : remplacer \ImagePresenceHPS par le chemin de l'image réelle.
    \safeincludegraphics[width=0.85\textwidth]{\ImagePresenceHPS}{Présence internationale de HPS}
    \caption{Présence internationale de HPS}
    \label{fig:presence-hps}
\end{figure}

\section{Historique de HPS}

Depuis sa création, HPS a progressivement développé son expertise autour de la conception de solutions logicielles dédiées aux paiements électroniques. Son évolution s'est construite autour de plusieurs axes : l'élargissement de son portefeuille de produits, l'ouverture vers les marchés internationaux, l'amélioration continue de ses processus internes et l'adoption de référentiels de qualité, de sécurité et de conformité adaptés aux exigences du secteur financier.

\begin{itemize}
    \item \textbf{1995} : création de HPS à Casablanca.
    \item \textbf{1996} : signature des premiers contrats.
    \item \textbf{2003} : ouverture de HPS Dubai et obtention de la certification ISO 9001.
    \item \textbf{2006} : introduction de HPS à la Bourse de Casablanca.
    \item \textbf{2010} : acquisition d'ACPQualife.
    \item \textbf{2012} : obtention d'un label relatif à la responsabilité sociale et environnementale.
    \item \textbf{2015} : création de la fondation HPS.
    \item \textbf{2016} : lancement de l'activité Processing pour le traitement des transactions de paiement.
    \item \textbf{2021} : acquisition d'ICPS et d'IPRC.
    \item \textbf{2024} : consolidation de la présence internationale de HPS dans plus de 90 pays.
\end{itemize}

\begin{figure}[h!]
    \centering
    % TODO : remplacer \ImageHistoriqueHPS par une frise chronologique réelle.
    \safeincludegraphics[width=0.85\textwidth]{\ImageHistoriqueHPS}{Historique de HPS}
    \caption{Historique de HPS}
    \label{fig:historique-hps}
\end{figure}

\section{Fiche technique de l'entreprise}

\begin{table}[h!]
    \centering
    \caption{Fiche technique de HPS}
    \label{tab:fiche-technique-hps}
    \begin{tabularx}{\textwidth}{>{\bfseries}p{4.2cm}X}
        \toprule
        Raison sociale & Hightech Payment Systems \\
        \midrule
        Sigle & HPS \\
        Secteur d'activité & Monétique et solutions de paiement électronique \\
        Année de création & 1995 \\
        Siège social & Casablanca, Maroc \\
        Présence internationale & Plus de 90 pays sur les cinq continents \\
        Effectif & Plus de 1000 experts et consultants \\
        Domaine d'expertise & Développement de solutions de paiement, traitement des transactions, sécurité logicielle et conformité \\
        Produit principal & Suite logicielle PowerCARD \\
        Marché cible & Banques, institutions financières, processeurs de paiement, switches nationaux et régionaux \\
        \bottomrule
    \end{tabularx}
\end{table}

\section{Organisation interne de HPS}

L'organisation interne de HPS repose sur une répartition fonctionnelle des responsabilités. Cette structuration vise à assurer la coordination entre les pôles métiers, les équipes techniques, les fonctions support et les instances de pilotage. Elle permet également d'aligner les objectifs opérationnels avec les exigences de qualité, de sécurité et de conformité propres au secteur des paiements.

De manière générale, l'entreprise s'organise autour de directions métiers, d'entités techniques, de fonctions de support et d'équipes projets. La Direction Software Factory occupe une place importante dans cette organisation, car elle porte l'industrialisation du développement logiciel et la structuration du cycle de vie des solutions.

\begin{figure}[h!]
    \centering
    % TODO : remplacer \ImageOrganigrammeHPS par l'organigramme officiel.
    \safeincludegraphics[width=0.9\textwidth]{\ImageOrganigrammeHPS}{Organigramme de HPS}
    \caption{Organigramme de HPS}
    \label{fig:organigramme-hps}
\end{figure}

\section{Direction Software Factory}

La Direction Software Factory a pour mission d'industrialiser le processus de développement logiciel. Elle vise à standardiser, automatiser et améliorer les différentes étapes du cycle de vie logiciel afin de produire des solutions robustes, maintenables et conformes aux exigences internes et externes.

Cette direction regroupe plusieurs branches liées à l'architecture interne, à l'architecture logicielle, au cycle de vie du développement logiciel et au développement de solutions. Dans le cadre d'une démarche DevSecOps, elle joue un rôle central dans la diffusion des pratiques de sécurité auprès des équipes projet.

\begin{figure}[h!]
    \centering
    % TODO : remplacer \ImageSoftwareFactory par le schéma réel.
    \safeincludegraphics[width=0.85\textwidth]{\ImageSoftwareFactory}{Direction Software Factory}
    \caption{Direction Software Factory}
    \label{fig:software-factory}
\end{figure}

\subsection{Equipe Software Security}

L'équipe d'accueil est l'équipe \textbf{Software Security}, rattachée au département SDLC de la Direction Software Factory. Elle intervient sur les sujets liés à la sécurité applicative, à l'industrialisation des contrôles et à l'amélioration continue du cycle de développement. Elle joue un rôle important dans la promotion des bonnes pratiques de développement sécurisé et dans l'accompagnement des équipes projet.

Ses missions couvrent notamment :

\begin{itemize}
    \item la définition et l'application des politiques de sécurité logicielle ;
    \item la mise en oeuvre des plans d'action issus de la qualité, des audits et des évaluations de sécurité ;
    \item l'installation, la configuration et la maintenance d'outils de sécurité ;
    \item la supervision des tests de pénétration et le suivi des vulnérabilités ;
    \item la gestion des incidents de sécurité logicielle et le suivi des correctifs ;
    \item la production de documentation liée au durcissement, à la modélisation des menaces et aux exigences de conformité ;
    \item la préparation et le maintien des certifications PCI SSS et PCI SSF ;
    \item la sensibilisation et la formation des développeurs aux pratiques de développement sécurisé.
\end{itemize}

\begin{figure}[h!]
    \centering
    % TODO : remplacer \ImageEquipeSecurity par le schéma de l'équipe.
    \safeincludegraphics[width=0.85\textwidth]{\ImageEquipeSecurity}{Equipe Software Security}
    \caption{Equipe Software Security}
    \label{fig:equipe-software-security}
\end{figure}

\subsection{Méthodologie de gestion de projet : Agile}

L'équipe Software Security adopte une méthodologie Agile favorisant la collaboration, la communication continue, l'adaptation aux changements, la livraison incrémentale et l'amélioration continue. Cette approche permet d'intégrer progressivement les exigences de sécurité dans les activités quotidiennes des équipes de développement.

Dans ce cadre, les échanges réguliers avec les équipes projet permettent d'identifier les besoins de sécurité, de suivre l'avancement des actions, de prioriser les corrections et de maintenir une vision partagée des risques applicatifs.

\begin{figure}[h!]
    \centering
    % TODO : remplacer \ImageAgile par une image de la méthodologie Agile utilisée.
    \safeincludegraphics[width=0.75\textwidth]{\ImageAgile}{Méthodologie Agile}
    \caption{Méthodologie de gestion de projet Agile}
    \label{fig:agile}
\end{figure}

\subsection{Outils de gestion de projet utilisés}

Les principaux outils utilisés dans ce cadre sont :

\begin{itemize}
    \item \textbf{Jira} : outil de gestion de projet utilisé pour planifier, suivre et piloter les tâches ;
    \item \textbf{Bitbucket / Git} : plateforme de gestion du code source et de suivi des versions ;
    \item \textbf{Confluence} : outil collaboratif utilisé pour documenter les processus, les décisions techniques et les bonnes pratiques.
\end{itemize}

\subsection{Implication d'un pentester dans la méthodologie Agile}

Dans une organisation Agile, le rôle du pentester consiste à contribuer à la sécurité du produit tout au long du cycle de développement. Ses responsabilités incluent la réalisation de tests manuels, l'élaboration de scénarios de test, la production de scripts d'automatisation, la participation aux réunions de suivi et la formulation de recommandations techniques à destination des équipes de développement.

Cette implication permet d'éviter que la sécurité soit traitée uniquement en fin de projet. Elle favorise au contraire une intégration continue des contrôles, des retours et des corrections, conformément aux principes DevSecOps.

\section{Solutions proposées par HPS}

HPS propose plusieurs solutions spécialisées dans le domaine de la monétique et des paiements électroniques. Ces solutions couvrent l'ensemble de la chaîne de valeur des paiements, depuis l'émission des cartes jusqu'au traitement des transactions, en passant par la gestion de la fraude, l'acquisition commerçant, la tokenisation et les services mobiles.

Parmi les principales solutions de la suite PowerCARD, on peut citer :

\begin{itemize}
    \item \textbf{PowerCARD-Issuer} : solution dédiée à l'émission et à la gestion des cartes ;
    \item \textbf{PowerCARD-Switch} : solution de traitement et de routage des transactions ;
    \item \textbf{PowerCARD-Fraud} : solution de surveillance et de détection de la fraude ;
    \item \textbf{PowerCARD-ATM} : solution de gestion des transactions au niveau des guichets automatiques ;
    \item \textbf{PowerCARD-Acquirer} : plateforme de gestion de l'acquisition commerçant ;
    \item \textbf{PowerCARD-WebPublisher} : solution de portail web pour les émetteurs et acquéreurs ;
    \item \textbf{PowerCARD-eSecure} : solution visant à renforcer la sécurité des paiements en ligne ;
    \item \textbf{PowerCARD-Wallet} : solution de paiement mobile ;
    \item \textbf{PowerCARD-Tokenisation} : solution de tokenisation et de détokénisation ;
    \item \textbf{PowerCARD-ACH} : solution de traitement et de règlement des transactions électroniques.
\end{itemize}

\section{Importance de la monétique chez HPS}

La monétique représente le coeur de l'activité de HPS. Elle regroupe l'ensemble des traitements électroniques, informatiques et télécoms nécessaires à la gestion des moyens de paiement et des transactions associées. Ce domaine couvre notamment les cartes bancaires, les paiements par internet, les paiements mobiles, les terminaux de paiement et les échanges entre systèmes financiers.

Dans ce contexte, HPS conçoit des solutions permettant aux institutions financières de gérer des transactions sensibles de manière sécurisée, fiable et conforme aux exigences du secteur. La sécurité logicielle, la traçabilité, la disponibilité et la conformité constituent donc des dimensions essentielles de ses activités.

\section{Bilan}

Ce chapitre a présenté Hightech Payment Systems, son historique, ses activités, son organisation interne ainsi que la Direction Software Factory et l'équipe Software Security. Il a également mis en évidence l'importance de la monétique dans les activités de HPS et le rôle central des pratiques de sécurité logicielle dans un environnement soumis à des exigences fortes de conformité.

Ces éléments permettent de mieux comprendre le contexte professionnel dans lequel s'inscrit le projet de fin d'études. Le chapitre suivant présente le cadrage du projet, sa problématique, ses objectifs et la méthodologie adoptée.

\clearpage
